\chapter{Herramientas empleadas}
\label{sec:herramientas}

La generación del conjunto de datos y la simulación del ruido se realizan con \textbf{Qiskit}
y su simulador Aer \cite{javadi2024qiskit}. Los modelos tabulares se implementan con
\textbf{scikit-learn} \cite{pedregosa2011scikit}, y la red neuronal de grafos con
\textbf{PyTorch} \cite{paszke2019pytorch} y \textbf{PyTorch Geometric}
\cite{fey2019pyg}. El seguimiento de los experimentos, con su semilla, la versión del código
y la huella del conjunto de datos, se lleva con \textbf{Weights \& Biases}
\cite{biewald2020wandb}.

El cómputo se reparte entre la máquina de desarrollo y \textbf{Google Cloud}
\cite{google2026cloud}. Esta última es una necesidad: las muestras de
mayor tamaño exigen más memoria de la que tiene la máquina local, y el entrenamiento de la
red de grafos requiere una unidad de procesamiento gráfico.
